% Drawing 147/03-5 -- folder 147, batch 03, archive page 53 (PDF page 54).
% Transcription: Les trois graphes de (65) sont dessinés : un cercle portant un sommet marqué « 1 » ; deux sommets marqués « 1 » reliés par deux arcs ; un cercle avec une queue dont le bout est marqué « 2 ». En regard, les courbes correspondantes : une courbe nodale (en huit) portant un point $a$ ; deux droites se coupant en deux points, portant $a$ et $b$ ; une courbe nodale à laquelle est attachée une droite portant $a$ et $b$.
% Kind: MD graphs of (65) and the corresponding stable curves. Method 3 (encoded, rendered from the data). Status: ok.
% Notes: Left: the three graphs of (65): a loop on a vertex of weight 1; two vertices of weight 1 joined by two arcs; a loop with a tail whose end has weight 2. Right: the corresponding curves: a nodal curve (figure eight) with a; two lines meeting in two points, with a and b; a nodal curve with a line attached carrying a and b. Encoded twice, the graphs (weights as legs) and the dual graphs of the curves: the two encodings coincide and give types (1,1), (1,2), (1,2) with ambient dimension 3g - 3 + nu = nu = 1, 2, 2, as in (65). The legends nu = ... = dim amb G are text. The labels a, b on the curves are small; the a of the first curve is not clearly visible.
% Coordinates: 1 unit = 1 mm = 1/100 of the crop width (bbox in page fractions: [0.347, 0.614, 0.9, 0.826]).
% Redrawn by Claude (claude-opus-5-5), issue #73, 2026-11-07; not checked by a human; not compiled (no TeX in the session).
% Rights: reproduces part of the work -- behind RENDER_TRANSCRIPTS, never published by itself.
\documentclass[tikz,border=3pt]{standalone}
\usepackage{amsmath,amssymb}
\usetikzlibrary{patterns,arrows.meta}
\tikzset{pencil/.style={black!45}, struck/.style={black!35}, mark/.style={circle,draw,inner sep=1.3pt}, pt/.style={circle,fill,inner sep=0.9pt}, opt/.style={circle,draw,fill=white,inner sep=1.1pt}}
\begin{document}
\begin{tikzpicture}[x=1mm,y=1mm,line cap=round,line join=round,font=\scriptsize]
\draw (18.8,-12.9) ellipse[x radius=4.7, y radius=3.5];
\node[pt] at (14.1,-12.9) {};
\node[inner sep=1pt] at (8.2,-10.6) {\tiny $1$};
\draw plot[smooth, tension=0.55] coordinates {(9.4,-22.4) (15.3,-19.4) (21.2,-22.4)};
\draw plot[smooth, tension=0.55] coordinates {(9.4,-22.4) (15.3,-25.3) (21.2,-22.4)};
\node[pt] at (9.4,-22.4) {};
\node[pt] at (21.2,-22.4) {};
\node[inner sep=1pt] at (7.1,-20) {\tiny $1$};
\node[inner sep=1pt] at (23.5,-20) {\tiny $1$};
\draw (18.8,-34.1) circle[radius=4.7];
\draw (14.1,-35.3) -- (5.9,-35.3);
\node[pt] at (14.1,-35.3) {};
\node[pt] at (5.9,-35.3) {};
\node[inner sep=1pt] at (5.9,-31.8) {\tiny $2$};
\draw plot[smooth, tension=0.55] coordinates {(72.9,-4.7) (78.8,-9.4) (83.5,-8.2) (89.4,-4.7) (94.1,-7.1) (89.4,-10.6) (83.5,-8.2) (78.8,-5.9) (72.9,-9.4)};
\node[inner sep=1pt] at (95.3,-3.5) {\tiny $a$};
\draw plot[smooth, tension=0.55] coordinates {(74.1,-21.2) (83.5,-14.1) (94.1,-12.9) (98.8,-12.9)};
\draw plot[smooth, tension=0.55] coordinates {(76.5,-15.3) (85.9,-22.4) (97.6,-20)};
\node[pt] at (80,-17.1) {};
\node[pt] at (91.8,-20) {};
\node[inner sep=1pt] at (77.6,-12.9) {\tiny $a$};
\node[inner sep=1pt] at (92.9,-23.5) {\tiny $b$};
\draw plot[smooth, tension=0.55] coordinates {(85.9,-22.4) (80,-30.6) (76.5,-35.3) (78.8,-37.6) (82.4,-31.8) (85.9,-27.1)};
\draw (68.2,-49.4) -- (84.7,-25.9);
\node[pt] at (72.9,-42.9) {};
\node[pt] at (70.6,-45.9) {};
\node[inner sep=1pt] at (75.3,-43.5) {\tiny $a$};
\node[inner sep=1pt] at (72.9,-48.2) {\tiny $b$};
\end{tikzpicture}
\end{document}
